\origpage{334}{N. Robertson, P.D. Seymour / Journal of Combinatorial Theory, Series B 92 (2004) 325--357}
\noindent where in the second case, $H$ is the rooted hypergraph such that $H \subseteq G_2$, $\pi(H) = \mu_2(e)$, $E(H) = \{e\}$ and $V(H)$ is the set of ends of $e$ in $G_2$. \fix{Since $\mathcal{F}_1$ is well-behaved, there exists $K_0$ such that $|\bar{\pi}(G)| \le K_0$ for every $(G, \mu, \Delta) \in \mathcal{F}_1$, and so $|\bar{\pi}(G_2)| \le K_0 + 1$. Thus we may split $\mathcal{F}_2$ into $K_0 + 2$ sets, according to whether $v \in \bar{\pi}(G_2)$ and, if so, which term of $\pi(G_2)$ it is. A union of finitely many well-behaved sets is well-behaved, so it suffices to show that each of these sets is well-behaved; and so we may assume that all members of $\mathcal{F}_2$ are in the same set.} Let us define $\gamma(Q_2) = Q_1$; then it is easy to see that $\gamma$ is an encoding and the result follows from 4.2.\fixnote{fix:4.9-proof}{As printed, $Q_1$ does not record whether $v$ is a root of $G_2$, or which term of $\pi(G_2)$ it is, but an expansion of $Q_2$ must send the $i$th root of $G_2$ to the $i$th root of $G'_2$. So $\gamma$ need not be an encoding. For example, let $\theta=1$, and let $\mathcal{F}_1$ consist of one patchwork, with one vertex $u$, root march $(u)$ and no edges. Let $G_2$ and $G'_2$ both have vertices $u, v$ and no edges, with $\pi(G_2)=(u,v)$ and $\pi(G'_2)=(u)$. Then $Q_2/\{v\}$ and $Q'_2/\{v\}$ are the same patchwork, so $\gamma(Q_2)$ is simulated in $\gamma(Q'_2)$. But $Q_2$ is not simulated in $Q'_2$, because an expansion keeps the length of the root march. The inserted split fixes the place of $v$, and then $\gamma$ is an encoding. The bound $K_0$ exists because expansions keep the number of roots. So 4.9 holds as printed.} $\square$

Let $P_1 = (G_1, \mu_1, \Delta_1)$ and $P_2 = (G_2, \mu_2, \Delta_2)$ be patchworks. We say that $P_1$ is a \textit{condensation} of $P_2$ if $V(G_1) = V(G_2)$, $\pi(G_1) = \pi(G_2)$, $dom(\mu_1) = E(G_1)$, $dom(\mu_2) = E(G_2)$, for each $e \in E(G_1)$ there is a rooted subhypergraph $A_e \subseteq G_2$ with the following properties:

\hspace*{12pt}$\bullet$ $V(A_e)$ is the set of ends of $e$ in $G_1$, and $\pi(A_e) = \mu_1(e)$,

\hspace*{12pt}$\bullet$ $\bigcup_{e \in E(G_1)} E(A_e) = E(G_2)$,

\hspace*{12pt}$\bullet$ for distinct $e, e' \in E(G_1)$, $E(A_e) \cap E(A_{e'}) = \emptyset$,

\hspace*{12pt}$\bullet$ for each $e \in E(G_1)$ and $K \in \Delta_1(e)$, $K$ is feasible in $P_2|A_e$.

A patchwork $P = (G, \mu, \Delta)$ is \textit{removable} if for every $e \in E(G)$, $\Delta(e)$ contains $N_V$ where $V$ is the set of ends of $e$.

\textbf{4.10.} \textit{Let $\mathcal{F}_1$ be a well-behaved set of removable patchworks and let $\mathcal{F}_2$ be a set of patchworks such that for each $P_2 \in \mathcal{F}_2$ some $P_1 \in \mathcal{F}_1$ is a condensation of $P_2$. Then $\mathcal{F}_2$ is well-behaved.}

\textbf{Proof.} Choose $N \ge 0$ (by 4.1) such that for every $(G, \mu, \Delta) \in \mathcal{F}_1$ and every $e \in dom(\mu)$, $e$ has $\le N$ ends. Now let $\Omega_2$ be a well-quasi-order. Let $\Omega_1$ be the well-quasi-order with $E(\Omega_1)$ the set of all $\Omega_2$-patchworks $(G, \mu, \Delta, \phi)$ where $|V(G)| \le N$, and $(G, \mu, \Delta)$ is removable, ordered by simulation (That $\Omega_1$ is a well-quasi-order \del{is proved in the same way as theorem 8.4 of [1] and we omit the proof} \fix{follows from theorem 8.4 of [1], applied to stars with one vertex. Simulation is a quasi-order, by theorem 7.1 of [1]. Let $P_i = (G_i, \mu_i, \Delta_i, \phi_i)$ $(i = 1, 2, \dots)$ be a countable sequence of members of $E(\Omega_1)$. For each $i \ge 1$ let $S_i$ be the star with one vertex $o(S_i)$ and no edges, and let $\sigma_i(o(S_i)) = G_i$. Then $(S_i, \sigma_i)$ is a star-decomposition of $G_i$: conditions (i), (ii) and (v) in the definition of a tree-decomposition in [1] hold since $\sigma_i(o(S_i)) = G_i$, condition (iv) holds since $t_1 = t_2 = t_3$ there, and conditions (iii) and (vi) hold vacuously, since they are about two distinct vertices or about an edge of $S_i$. Also $|V(\sigma_i(o(S_i)))| = |V(G_i)| \le N$ and $P_i$ is removable, since $P_i \in E(\Omega_1)$. Since $U(S_i) = \emptyset$, the set $\{P_i | \sigma_i(s) \colon i \ge 1, s \in U(S_i)\}$ is empty, and so it is well-quasi-ordered by simulation. Hence by theorem 8.4 of [1], with $\Omega = \Omega_2$ and $n = N$, there exist $i' > i \ge 1$ such that $P_i$ is simulated in $P_{i'}$; and so $\Omega_1$ is a well-quasi-order}.\fixnote{fix:4.10-omega-proof-added}{The paper omits the proof that $\Omega_1$ is a well-quasi-order. The argument above is the one used in the formalization.})

Now let $Q_2 = (G_2, \mu_2, \Delta_2, \phi_2)$ be an $\Omega_2$-completion of some $P_2 \in \mathcal{F}_2$. Choose $P_1 = (G_1, \mu_1, \Delta_1) \in \mathcal{F}_1$ such that $P_1$ is a condensation of $P_2$, and choose the rooted subhypergraphs $A_e\ (e \in E(G_1))$ as in the definition of condensation. Let $Q_1 = (G_1, \mu_1, \Delta_1, \phi_1)$ be the $\Omega_1$-completion of $P_1$ where $\phi_1(e) = Q_2|A_e$ for each $e \in dom(\mu_1) = E(G_1)$. Let $Q_1 = \gamma(Q_2)$; then \del{theorem 5.7 of [4] implies} \fix{one can check directly} that $\gamma$ is an encoding, and the result follows.\fixnote{fix:4.10-proof}{Theorem 5.7 of [4] does not apply as stated. It needs one heart edge for each member of $\{A_e \colon e \in E(G_1)\}$, but two parallel edges $e \ne e'$ of $G_1$ with the same march and $E(A_e) = E(A_{e'}) = \emptyset$ have $A_e = A_{e'}$. Also, the patches $\Delta_1(e)$ are only subsets of $gr(P_2|A_e)$. The direct check uses that every vertex of each $A_e$ is a root of $A_e$. An expansion $\eta$ of $Q_1$ in $Q'_1$ gives an expansion of $Q_2$ in $Q'_2$ that sends each edge $f \in E(A_e)$ to the matching edge of $A'_{\eta(e)}$, and a realization of it is put together from realizations of the pieces that $\eta$ misses. So 4.10 holds as printed.} $\square$

\textbf{4.11.} \textit{Let $\mathcal{F}_1$ be a well-behaved set of active patchworks, and let $\mathcal{F}_2$ be the set of all patchworks $P_2 = (G_2, \mu_2, \Delta_2)$ such that there exists $P_1 = (G_1, \mu_1, \Delta_1) \in \mathcal{F}_1$ with $G_2 \subseteq G_1$, $\pi(G_2) = \pi(G_1)$, $G_2$ complemented in $G_1$ and $P_2 = P_1|G_2$. Then $\mathcal{F}_2$ is well-behaved.}

\textbf{Proof.} Let $\Omega_2$ be a well-quasi-order. Let $* \notin E(\Omega_2)$ be a new element and let $\Omega_1$ be the well-quasi-order with $\Omega_2 \subseteq \Omega_1$ and $E(\Omega_1) = E(\Omega_2) \cup \{*\}$, where if $x \le *$ or $* \le x$ then $x = *$. Now let $Q_2 = (G_2, \mu_2, \Delta_2, \phi_2)$ be an $\Omega_2$-completion of $P_2 = (G_2, \mu_2, \Delta_2) \in \mathcal{F}_2$. Choose $P_1 = (G_1, \mu_1, \Delta_1) \in \mathcal{F}_1$ so that $G_2 \subseteq G_1$, $\pi(G_2) = \pi(G_1)$, $G_2$ is complemented in $G_1$, and $P_2 = P_1|G_2$. Let $Q_1 = (G_1, \mu_1, \Delta_1, \phi_1)$ be the $\Omega$-completion of $P_1$ where
\begin{align*}
\phi_1(e) &= \phi_2(e)\ (e \in E(G_2)) \\
&= *\ (e \in E(G_1) \setminus E(G_2)).
\end{align*}
Let $\gamma(Q_2) = Q_1$; we claim that $\gamma$ is an encoding.
